\chapter{紧致度量空间}
\label{chap:compact-metric-spaces}

在欧氏空间中，紧致、序列紧致和闭有界等价；证明同时使用了有限维坐标与
有限网。一般度量空间仍有开球和数列，却未必有坐标。要恢复同样的紧致性
判据，需要把欧氏有界集的“每个尺度都有有限网”单独提出，这就是全有界性。
紧致度量空间恰好是既完备又全有界的空间。

\section{开覆盖紧致和序列紧致}

\begin{definition}[紧致与序列紧致]
度量空间 \(X\) 称为紧致的，如果每个开覆盖都有有限子覆盖。若 \(X\) 中
每个数列都有收敛到 \(X\) 中某点的子列，则称 \(X\) 序列紧致。
\end{definition}

\begin{lemma}[紧致空间中的无限集]
\label{lem:compact-metric-infinite-set-cluster}
紧致度量空间中的每个无限子集都有聚点。
\end{lemma}

\begin{proof}
反设无限集 \(A\subset X\) 没有聚点。若 \(x\in A\)，可取开邻域
\(U_x\) 使 \(U_x\cap A=\{x\}\)；若 \(x\notin A\)，可取 \(U_x\) 与
\(A\) 不相交。开集族 \(\{U_x:x\in X\}\) 覆盖 \(X\)，紧致性给出有限
子覆盖。但每个被选出的 \(U_x\) 至多含 \(A\) 的一个点，有限个这样的集合
不可能覆盖无限集 \(A\)，矛盾。
\end{proof}

\begin{theorem}[度量空间中两种紧致性的等价]
\label{thm:metric-compact-sequential-equivalence}
度量空间紧致，当且仅当它序列紧致。
\end{theorem}

\begin{proof}
空空间同时满足两种紧致性，以下设 \(X\ne\varnothing\)。
先设 \(X\) 紧。给定数列 \((x_n)\)。若某个值出现无穷多次，便有常值
子列。否则值集必为无限集，从原列中按首次出现的次序抽出一条各项互异的
子列，并重新记作 \((x_n)\)。其值集有聚点 \(x\)。每个
\(B(x,1/k)\) 含值集中的无穷多个点，否则缩小半径即可避开
其中除 \(x\) 外的有限多个点，与聚点定义矛盾。于是可递归选择
\[
 n_1<n_2<\cdots,\qquad x_{n_k}\in B(x,1/k),
\]
得到 \(x_{n_k}\to x\)。

反过来，设 \(X\) 序列紧致。先证明两个辅助事实。

第一，对每个 \(\varepsilon>0\)，有限多个半径为 \(\varepsilon\) 的开球
可以覆盖 \(X\)。否则可递归取点 \(x_n\)，使
\[
 d(x_n,x_j)\ge\varepsilon\qquad(j<n).
\]
这条序列的任意子列都不是 Cauchy 列，因而不可能收敛，与序列紧致矛盾。

第二，\(X\) 的每个开覆盖 \(\mathcal U\) 都有 Lebesgue 数：存在
\(\lambda>0\)，使每个 \(B(x,\lambda)\) 都包含于某个 \(U\in\mathcal U\)。
若不存在，可对每个 \(n\) 取 \(x_n\)，使
\(B(x_n,1/n)\) 不包含于任何覆盖成员。抽取
\(x_{n_k}\to x\)。选 \(U\in\mathcal U\) 含 \(x\)，再取
\(r>0\) 使 \(B(x,r)\subset U\)。充分大时
\(d(x_{n_k},x)<r/2\) 且 \(1/n_k<r/2\)，于是
\(B(x_{n_k},1/n_k)\subset U\)，矛盾。

取有限 \(\lambda/2\)-球覆盖，中心记为 \(a_1,\ldots,a_N\)。为每个中心
选择 \(U_j\in\mathcal U\)，使 \(B(a_j,\lambda)\subset U_j\)。因
\(B(a_j,\lambda/2)\subset U_j\)，有限多个 \(U_j\) 覆盖 \(X\)。所以
\(X\) 紧。
\end{proof}

等价证明依赖度量提供的可数尺度 \(1/n\) 和有限半径球。一般拓扑空间中，
开覆盖紧致与序列紧致可以分离。
\chapterref{chap:compact-topological-spaces}和
\chapterref{chap:nets-filters}会分别采用覆盖、网与滤子处理这一差别。

\section{全有界性}

非空集合 \(A\subset X\) 的直径定义为
\[
 \operatorname{diam}A=\sup\{d(x,y):x,y\in A\},
\]
允许取 \(+\infty\)，并约定 \(\operatorname{diam}\varnothing=0\)。
非空集合有界当且仅当直径有限；空集也视为有界集。

\begin{definition}[全有界]
度量空间 \(X\) 称为全有界的，如果对每个 \(\varepsilon>0\)，存在有限
\(\varepsilon\)-网，即有限点集 \(a_1,\ldots,a_N\in X\) 满足
\[
 X=\bigcup_{j=1}^N B(a_j,\varepsilon).
\]
\end{definition}

这里允许 \(N=0\)，空网覆盖空空间。非空全有界空间必有界：取
\(\varepsilon=1\) 的有限网和其中一个中心，用有限多个
中心间距离的最大值控制整个空间。反向不成立；无限离散度量空间的直径为一，
却在 \(\varepsilon<1\) 时每个球只有一个点，不可能有有限网。

\begin{definition}[\(\varepsilon\)-分离集]
集合 \(A\subset X\) 称为 \(\varepsilon\)-分离的，如果任意不同
\(x,y\in A\) 都满足 \(d(x,y)\ge\varepsilon\)。
\end{definition}

\begin{theorem}[全有界性的序列判据]
\label{thm:total-bounded-characterizations}
对度量空间 \(X\)，以下条件等价：
\begin{enumerate}[label=\textup{(\roman*)}]
 \item \(X\) 全有界；
 \item \(X\) 中每个数列都有 Cauchy 子列；
 \item 对每个 \(\varepsilon>0\)，\(X\) 中的
 \(\varepsilon\)-分离集都是有限集；
 \item 对每个 \(\varepsilon>0\)，\(X\) 可由有限多个直径小于
 \(\varepsilon\) 的集合覆盖。
\end{enumerate}
\end{theorem}

\begin{proof}
空空间满足四项条件，以下设空间非空。
(i) 推出 (ii)。依次取有限 \(1,1/2,1/3,\ldots\)-网。在第一个有限球
覆盖中，至少一个球含原序列的无穷多项；限制到这些项，在第二个覆盖中再取
含无穷多项的球，如此继续。第 \(k\) 步保留的指标集嵌套，并且相应各点都
落在同一半径 \(1/k\) 的球中。递归选择严格递增指标 \(n_k\)，使从第
\(k\) 项起都留在第 \(k\) 步的指标集中。若 \(p,q\ge k\)，则
\[
 d(x_{n_p},x_{n_q})<2/k,
\]
所以所得子列为 Cauchy 列。

(ii) 推出 (iii)。若存在无限 \(\varepsilon\)-分离集，从中选两两不同的
点列；它的任意子列中不同项距离都不小于 \(\varepsilon\)，不可能是
Cauchy 列。

(iii) 推出 (i)。固定 \(\varepsilon>0\)，从任一点开始逐次加入与已选
各点距离至少为 \(\varepsilon\) 的新点。若过程不停止，所选点构成无限
\(\varepsilon\)-分离集；所以有限步后停止。此时极大性说明每个 \(x\in X\)
都距某个已选点小于 \(\varepsilon\)，故这些点构成有限网。

(i) 推出 (iv)：用半径 \(\varepsilon/3\) 的有限球覆盖，每个球的直径至多为
\(2\varepsilon/3<\varepsilon\)。(iv) 推出 (i)：在每个非空覆盖成员中
选一点；若该集合直径小于 \(\varepsilon\)，它包含于以所选点为中心的
\(\varepsilon\)-球。有限多个所选点构成有限网。
\end{proof}

\begin{proposition}[全有界性的保持]
\label{prop:total-bounded-preservation}
若 \(A\) 是度量空间 \(Y\) 的全有界子集，则 \(A\) 的每个子集以及
它在 \(Y\) 中的闭包都全有界；有限多个全有界子集的并全有界。一致连续
映射把全有界空间映为全有界空间。
\end{proposition}

\begin{proof}
子集可沿用原空间有限网，但若要求网点属于子集，就从每个与子集相交的网球
中另选一点，并把半径放大一倍。对闭包，先以半径 \(\varepsilon/2\) 覆盖
原集合；任一闭包点可由原集合点逼近到 \(\varepsilon/2\) 内，所以落在同一
中心的 \(\varepsilon\)-球中。有限并分别取有限网后合并。

设 \(f:X\to Y\) 一致连续。给定 \(\varepsilon>0\)，取 \(\delta>0\)
使 \(d_X(x,y)<\delta\) 推出
\(d_Y(f(x),f(y))<\varepsilon\)。取 \(X\) 的有限 \(\delta\)-网，网点的
像便构成 \(f(X)\) 的有限 \(\varepsilon\)-网。
\end{proof}

\section{紧致等价于完备且全有界}

\subsection{完备性补上极限}

\begin{theorem}[紧致度量空间的判据]
\label{thm:compact-iff-complete-totally-bounded}
度量空间 \(X\) 紧，当且仅当 \(X\) 完备且全有界。
\end{theorem}

\begin{proof}
若 \(X\) 紧，则由上一节证明中的有限球覆盖可知它全有界。任取 Cauchy 列，
序列紧致性给出收敛子列；Cauchy 列只要有一个收敛子列，整列就收敛到同一
极限，所以 \(X\) 完备。

反过来，若 \(X\) 全有界，则每个数列都有 Cauchy 子列；完备性使这条子列
在 \(X\) 中收敛。因此 \(X\) 序列紧致，再由度量空间中两种紧致性的等价
得到紧致性。
\end{proof}

两个条件承担不同任务。全有界性负责从任意数列中取得 Cauchy 子列，完备性
负责让这条子列在空间内有极限。只保留前者时，\((0,1)\) 全有界却不紧；
只保留后者时，无限离散度量空间完备却不紧。

\subsection{预紧性与完备化}

有些文献称闭包紧致的集合为相对紧或预紧。对度量空间，全有界正好描述
其完备化中的这种性质。

\begin{corollary}[全有界空间的完备化]
\label{cor:completion-total-bounded-compact}
度量空间 \(X\) 全有界，当且仅当它的完备化 \(\widehat X\) 紧。
\end{corollary}

\begin{proof}
若 \(X\) 全有界，它在 \(\widehat X\) 中的像全有界，闭包
\(\widehat X\) 也全有界；完备化本身完备，故紧。反之，若
\(\widehat X\) 紧，则它全有界，其子空间 \(X\) 也全有界。
\end{proof}

\begin{corollary}[有限乘积]
有限多个紧致度量空间配最大值度量的乘积仍紧。
\end{corollary}

\begin{proof}
有限乘积的 Cauchy 列逐坐标 Cauchy，所以完备性由各因子完备性得到。对
给定 \(\varepsilon\)，分别取各因子的有限 \(\varepsilon\)-网，所有网点的
有限笛卡儿积构成乘积空间的有限 \(\varepsilon\)-网。判据给出紧致性。
\end{proof}

\section{紧致空间上的一致连续性}

\begin{theorem}[度量空间的 Heine--Cantor 定理]
\label{thm:heine-cantor-metric}
若 \(X\) 紧，\(f:X\to Y\) 连续，则 \(f\) 一致连续。
\end{theorem}

\begin{proof}
若不一致连续，存在 \(\varepsilon_0>0\) 及点列 \(x_n,y_n\in X\)，使
\[
 d_X(x_n,y_n)<1/n,\qquad
 d_Y(f(x_n),f(y_n))\ge\varepsilon_0.
\]
由序列紧致性，某子列 \(x_{n_k}\to x\in X\)。三角不等式给出
\(y_{n_k}\to x\)。连续性迫使两条像子列都趋于 \(f(x)\)，故二者距离趋零，
与统一下界矛盾。
\end{proof}

紧致性还把逐点选取的邻域尺度压缩成一个覆盖全空间的尺度。这个结论在
分割区间、拼接局部估计和证明同伦性质时都会使用。

\begin{definition}[Lebesgue 数]
设 \(\mathcal U\) 是度量空间 \(X\) 的一个开覆盖。若 \(\lambda>0\)
满足：对每个 \(x\in X\)，都存在 \(U\in\mathcal U\) 使
\[
 B(x,\lambda)\subseteq U,
\]
则称 \(\lambda\) 为覆盖 \(\mathcal U\) 的一个\term{Lebesgue 数}。
\end{definition}

\begin{theorem}[Lebesgue 数引理]
\label{thm:lebesgue-number-compact-metric}
紧致度量空间的每个开覆盖都有 Lebesgue 数。
\end{theorem}

\begin{proof}
设开覆盖 \(\mathcal U\) 没有 Lebesgue 数。对每个 \(n\in\Npos\)，
可取 \(x_n\in X\)，使 \(B(x_n,1/n)\) 不包含于 \(\mathcal U\) 的任何
成员。紧致性给出收敛子列 \(x_{n_k}\to x\)。选取
\(U_0\in\mathcal U\) 使 \(x\in U_0\)，再取 \(r>0\) 使
\(B(x,r)\subseteq U_0\)。当 \(k\) 充分大时，
\[
 d(x_{n_k},x)<r/2,
 \qquad
 1/n_k<r/2.
\]
于是三角不等式给出
\(B(x_{n_k},1/n_k)\subseteq B(x,r)\subseteq U_0\)，与
\(x_{n_k}\) 的选取矛盾。
\end{proof}

定义中使用球并非唯一选择。若每个直径小于 \(\lambda\) 的非空子集都包含于
某个覆盖成员，则每个半径小于 \(\lambda/2\) 的球都被某个成员包含；
反过来，在非空子集中任选一点作为球心，球版本便控制其余各点。
两种表述只相差一个
固定常数。

因此紧空间上的连续映射保持 Cauchy 列。更一般地，一致连续映射把全有界
空间映成全有界空间；给定像空间中的误差，只需先选取相应的一致连续尺度
\(\delta\)，再对定义域取有限 \(\delta\)-网。

欧氏空间中紧集的几个保持性质在度量空间中仍成立，证明如下。
若 \(F\) 是紧空间 \(X\) 的闭子集，把 \(F\) 的相对开覆盖写成
\(F\cap U_\alpha\)，其中 \(U_\alpha\) 在 \(X\) 中开，再加入
\(X\setminus F\)，便得到 \(X\) 的开覆盖。取有限子覆盖后与 \(F\) 相交，
即知 \(F\) 紧。若 \(f:X\to Y\) 连续，把 \(f(X)\) 的相对开覆盖取原像，
在 \(X\) 中选有限子覆盖再映回，便知 \(f(X)\) 紧。

紧子集 \(K\subset Y\) 必闭。若 \(y\in\overline K\)，闭包的序列刻画给出
\(x_n\in K\) 趋于 \(y\)；\(K\) 的序列紧致性给出子列趋于某个 \(x\in K\)。
极限唯一性迫使 \(y=x\in K\)。因此紧空间上的实值连续函数的像是实紧集，
必有界；若定义域非空，它还取得最大值和最小值。若
\(f:X\to Y\) 是紧空间到度量空间的连续双射，则对任意闭集
\(F\subset X\)，集合 \(F\) 紧，\(f(F)\) 紧而在 \(Y\) 中闭，所以
\(f^{-1}\) 连续。由此得到紧空间上的连续双射自动为同胚。

紧致性不是一致连续的必要条件。以离散度量为定义域时，任意映射都一致连续：
给定任意正误差，只需取 \(\delta<1\)，条件 \(d_X(x,y)<\delta\) 就迫使
\(x=y\)。通常距离下的恒等映射
\(\R\to\R\) 也一致连续，而 \(\R\) 不紧。

\section{Cantor 交集定理及闭有界反例}

\subsection{直径趋零的闭集套}

\begin{theorem}[Cantor 交集定理]
\label{thm:cantor-intersection-complete-metric}
度量空间 \(X\) 完备，当且仅当每个满足
\[
 F_1\supset F_2\supset\cdots,\qquad
 F_n\ne\varnothing,\quad F_n\text{ 闭},\quad
 \operatorname{diam}F_n\to0
\]
的闭集套都有非空交。此时交集恰含一个点。
\end{theorem}

\begin{proof}
先设 \(X\) 完备。从每个 \(F_n\) 取 \(x_n\)。给定 \(\varepsilon>0\)，
取 \(N\) 使 \(\operatorname{diam}F_N<\varepsilon\)。若 \(m,n\ge N\)，
则 \(x_m,x_n\in F_N\)，所以 \(d(x_m,x_n)<\varepsilon\)。故
\((x_n)\) 为 Cauchy 列，收敛到某个 \(x\in X\)。固定 \(N\) 后，尾列
都在闭集 \(F_N\) 中，故 \(x\in F_N\)。于是
\(x\in\bigcap_nF_n\)。若 \(x,y\) 都属于交集，则
\(d(x,y)\le\operatorname{diam}F_n\) 对每个 \(n\) 成立，令 \(n\to\infty\)
得 \(x=y\)。

反过来，假设闭集套性质成立，令 \((x_n)\) 为任意 Cauchy 列。定义
\[
 F_n=\overline{\{x_k:k\ge n\}}.
\]
这些集合非空、闭且递减。给定 \(\varepsilon>0\)，Cauchy 性给出 \(N\)，
使 \(k,\ell\ge N\) 时 \(d(x_k,x_\ell)<\varepsilon/2\)。若 \(u,v\in F_N\)，
分别取尾集中的点列趋于 \(u,v\)。每对逼近点的距离小于 \(\varepsilon/2\)，
取极限得到 \(d(u,v)\le\varepsilon/2\)。故
\(\operatorname{diam}F_N\le\varepsilon/2<\varepsilon\)；于是直径趋零。
交集中存在唯一点 \(x\)。因 \(x\in F_n\) 且
\(\operatorname{diam}F_n\to0\)，有
\(d(x_n,x)\le\operatorname{diam}F_n\to0\)。所以每个 Cauchy 列收敛，
\(X\) 完备。
\end{proof}

\subsection{紧集套不需要直径条件}

\begin{theorem}[紧空间中的闭集套]
\label{thm:nested-closed-compact-metric}
若 \(X\) 紧，\(F_1\supset F_2\supset\cdots\) 是非空闭集套，则
\(\bigcap_nF_n\ne\varnothing\)；这里不要求直径趋零。
\end{theorem}

\begin{proof}
若交为空，则开集 \(X\setminus F_n\) 覆盖 \(X\)。取有限子覆盖并令 \(N\)
为其中最大指标。因集合递减，补集递增，所以
\[
 X=\bigcup_{n\le N}(X\setminus F_n)=X\setminus F_N,
\]
这迫使 \(F_N=\varnothing\)，与非空性矛盾。
\end{proof}

直径条件在完备空间中补足紧致性所提供的全有界控制；没有它，完备空间中的
闭集套可以有空交，例如在 \(\R\) 中取 \(F_n=[n,\infty)\)。

\subsection{闭有界不再充分}

无限集合配离散度量时，整个空间闭且有界，并且完备，但不全有界，因而不紧。
同样的现象出现在函数空间
\(\mathcal B(\Npos)\) 的闭单位球中：标准函数 \(e_n\) 两两距离为一，
所以闭单位球不可能被有限多个半径小于 \(1/2\) 的球覆盖。一般度量空间中，
“闭且有界”只保留了欧氏 Heine--Borel 定理的两个必要条件；充分条件应换成
“完备且全有界”。

\section*{习题}
\addcontentsline{toc}{section}{习题}

\subsection*{A组　紧致与全有界}
\begin{enumerate}[label=\textbf{\thechapter.\arabic*.}]
 \exerciseitem{compact-infinite-set-cluster}
 完整证明紧致度量空间中的无限集有聚点，并处理覆盖中心属于或不属于该
 无限集的两种情形。
 \exerciseitem{compact-to-sequential-subsequence}
 从聚点引理构造收敛子列，说明为什么聚点的每个邻域必须含原序列的无穷多项。
 \exerciseitem{sequential-compact-lebesgue-number}
 证明 Lebesgue 数的球版本与“小直径集合版本”等价到一个固定因子；再从
 序列紧致性证明任意开覆盖有 Lebesgue 数，并用有限球覆盖推出紧致性。
 \exerciseitem{total-bounded-implies-bounded}
 证明全有界空间有界，并给出有界但不全有界的完备空间。
 \exerciseitem{separated-set-characterization}
 证明全有界性等价于每个 \(\varepsilon\)-分离集有限。
 \exerciseitem{total-bounded-cauchy-subsequence}
 把有限 \(1/k\)-网的逐层抽取写成严格递增指标，证明所得子列为 Cauchy 列。
 \exerciseitem{total-bounded-closure}
 证明全有界子集的闭包仍全有界；核对为何需要把初始网半径缩小一半。
\end{enumerate}

\subsection*{B组　完备性与紧致性}
\begin{enumerate}[label=\textbf{\thechapter.\arabic*.},resume]
 \exerciseitem{compact-complete-proof}
 证明紧致度量空间完备，且不得直接引用“紧致等价于完备且全有界”的结论。
 \exerciseitem{complete-total-bounded-compact}
 从“每列有 Cauchy 子列”重建完备且全有界推出紧致的证明链。
 \exerciseitem{completion-compact-characterization}
 证明 \(X\) 全有界当且仅当其完备化紧。
 \exerciseitem{finite-product-total-bounded}
 证明有限多个全有界空间的最大值度量乘积全有界；说明为什么任意多个因子的
 结论不能直接沿用有限网的笛卡儿积。
 \exerciseitem{uniform-image-total-bounded}
 证明一致连续像保持全有界；构造连续但非一致连续映射把全有界空间映成
 非有界空间。
 \exerciseitem{metric-heine-cantor}
 分别用序列紧致性和 Lebesgue 数证明度量空间版 Heine--Cantor 定理。
\end{enumerate}

\subsection*{C组　闭集套与反例}
\begin{enumerate}[label=\textbf{\thechapter.\arabic*.},resume]
 \exerciseitem{cantor-intersection-forward}
 在完备空间中证明直径趋零的非空闭集套交于唯一一点，并逐项指出闭性、
 完备性与直径条件的使用位置。
 \exerciseitem{cantor-intersection-converse}
 从 Cauchy 列的尾集闭包构造闭集套，证明 Cantor 交集性质推出完备性。
 \exerciseitem{compact-nested-closed}
 用开覆盖定义证明紧空间中非空闭集套交非空，并给出交集含多个点的例子。
 \exerciseitem{closed-bounded-discrete-counterexample}
 在无限离散度量空间中验证整个空间闭、有界、完备但不紧。
 \exerciseitem{closed-ball-function-space-not-compact}
 在 \(\mathcal B(\Npos)\) 的闭单位球中构造无限 \(1\)-分离集，由此证明
 该闭球不紧。
 \optionalexerciseitem{compactness-finite-intersection-property}
 证明空间 \(X\) 紧，当且仅当每族具有有限交性质的闭集的总交非空。
 有限交性质指每个有限子族的交非空，包含空子族；空子族的交约定为 \(X\)。
\end{enumerate}
